\chapter{Hello, Lattice}\label{ch:hello-lattice}

\index{Lattice!introduction}

Every programming language begins with a promise.
Some promise speed.
Others promise safety, or simplicity, or raw expressive power.
Lattice promises something different: it promises you \emph{control over change itself}.

In most languages, the question of whether a value can be modified---or when it should stop being modified---is answered by convention, or by a single keyword like \lstinline|const|.
Lattice takes that question seriously and builds an entire system around it, drawing on a metaphor from chemistry and materials science.
Values in Lattice are like materials: they can be \emph{fluid}, flowing and reshaping freely, or they can \emph{crystallize} into a permanent, immutable form.
This is the \emph{phase system}\index{phase system}, and it is woven into every corner of the language.

But we are getting ahead of ourselves.
Before we explore crystallization, let's meet Lattice the way any programmer meets a new language---by writing our first program and watching it run.

%% ───────────────────────────────────────────────────────────
\section{What Lattice Is and Why It Exists}\label{sec:what-is-lattice}
\index{Lattice!design philosophy}

Lattice is an interpreted, general-purpose programming language implemented in C\@.
It compiles your source code to bytecode and executes it on a stack-based virtual machine, with the option to run on a register-based VM or a tree-walking interpreter as well.
The current release is v0.3.28, and it runs on macOS, Linux, and (experimentally) via WebAssembly in the browser.

On the surface, Lattice looks familiar.
If you've written Python, JavaScript, Rust, or Go, you'll recognize the curly braces, the \lstinline|if|/\lstinline|else| blocks, and the \lstinline|fn| keyword for functions:

\begin{lstlisting}[language=Lattice, caption={A first glance at Lattice syntax}]
fn greet(name: String) -> String {
    "Hello, ${name}!"
}

print(greet("world"))
// Hello, world!
\end{lstlisting}

So what makes Lattice worth learning?
Three things stand out.

\subsection{The Phase System}\label{subsec:phase-promise}
\index{phase system!overview}

Most languages draw a binary line: a variable is either mutable or immutable.
Lattice replaces that binary with a richer model.
A value starts \emph{fluid}\index{fluid phase}---you can modify it, grow it, reshape it.
When you are done, you \emph{freeze}\index{freeze} it, and it becomes \emph{crystal}\index{crystal phase}---locked, permanent, safe to share across threads.
If you need to modify it again later, you \emph{thaw}\index{thaw} it, which gives you a mutable copy without disturbing the original.

This is not merely a convention.
The Lattice runtime enforces it.
Try to mutate a frozen value and you will get an error---not a silent bug discovered in production at 2 a.m.

\begin{lstlisting}[language=Lattice, caption={Freezing and thawing}]
flux temperatures = [72.1, 68.5, 74.3]
temperatures.push(71.0)   // fine: temperatures is fluid

freeze(temperatures)
// temperatures.push(80.0) -- error! crystal values cannot be modified

flux updated = thaw(temperatures)  // thaw produces a fluid copy
updated.push(80.0)                 // now we can modify again
\end{lstlisting}

We will explore the phase system gradually throughout this book, with a full deep-dive in Part~III\@.

\subsection{Familiar Syntax, Modern Features}\label{subsec:modern-features}

Lattice ships with first-class closures\index{closures}, pattern matching\index{pattern matching} with exhaustiveness checking, structs with traits and \lstinline|impl| blocks, string interpolation\index{string interpolation}, lazy iterators\index{iterators}, structured concurrency\index{concurrency} with channels, and a built-in test framework.
It is a language designed for working programmers who want to get things done without drowning in ceremony.

\subsection{Batteries Included}\label{subsec:batteries}

Lattice includes over 120 built-in functions\index{built-in functions} spanning file I/O, HTTP, JSON/TOML/YAML parsing, regular expressions, cryptography, networking, and more.
The standard library modules can be imported with a single line.
There is also a package manager for third-party libraries.

\begin{notebox}{Lattice Is Young}
Lattice is at version 0.3.28.
The language is actively evolving.
Some features described in this book may change in future releases, though the core concepts---the phase system, the type model, the concurrency story---are stable foundations that the rest of the language builds upon.
\end{notebox}

%% ───────────────────────────────────────────────────────────
\section{Installing Lattice}\label{sec:installing-lattice}
\index{installation}

Lattice is built from source using a C11 compiler and \texttt{make}.
There is no pre-built binary distribution yet---part of the charm of a young language is that you get to compile it yourself.

\subsection{Prerequisites}\label{subsec:prerequisites}

You need two things:

\begin{enumerate}
  \item A \textbf{C11-compatible compiler}\index{C11 compiler}: GCC or Clang on Linux, Apple Clang on macOS, or MinGW on Windows.
  \item \textbf{libedit}\index{libedit}: a lightweight readline library used for the REPL's line editing and tab completion. It ships with macOS. On Linux, install the development package.
\end{enumerate}

Optionally, if you want TLS networking (\lstinline|tls_connect|, \lstinline|tls_read|, etc.) and cryptographic built-ins (\lstinline|sha256|, \lstinline|md5|, \lstinline|base64_encode|), you'll need \textbf{OpenSSL}\index{OpenSSL} available via \texttt{pkg-config}.

\subsection{macOS}\label{subsec:install-macos}
\index{installation!macOS}

macOS ships with Apple Clang and libedit.
Clone the repository and build:

\begin{lstlisting}[language=bash, numbers=none]
git clone https://github.com/alexjokela/lattice.git
cd lattice
make
\end{lstlisting}

This produces a binary called \texttt{clat}\index{clat} in the current directory.
You can move it to a directory on your \texttt{PATH}:

\begin{lstlisting}[language=bash, numbers=none]
sudo cp clat /usr/local/bin/
\end{lstlisting}

\subsection{Linux}\label{subsec:install-linux}
\index{installation!Linux}

On Debian or Ubuntu, install the libedit development headers first:

\begin{lstlisting}[language=bash, numbers=none]
sudo apt install libedit-dev
git clone https://github.com/alexjokela/lattice.git
cd lattice
make
\end{lstlisting}

On Fedora or RHEL:

\begin{lstlisting}[language=bash, numbers=none]
sudo dnf install libedit-devel
\end{lstlisting}

\subsection{Windows}\label{subsec:install-windows}
\index{installation!Windows}

Lattice includes a Win32 compatibility layer (\filename{src/win32\_compat.c}) that provides POSIX-like functions.
Using MinGW or WSL:

\begin{lstlisting}[language=bash, numbers=none]
# Using WSL (recommended)
sudo apt install libedit-dev build-essential
cd lattice
make
\end{lstlisting}

\begin{tipbox}{WSL Is the Smoothest Path}
If you are on Windows, the Windows Subsystem for Linux gives you the best experience.
Lattice was primarily developed on macOS and Linux, and WSL brings you into that environment without dual-booting.
\end{tipbox}

\subsection{WebAssembly}\label{subsec:install-wasm}
\index{installation!WebAssembly}\index{WebAssembly}

Lattice can be compiled to WebAssembly using Emscripten, allowing it to run in a browser.
This is experimental, but it opens up interesting possibilities for playgrounds and education:

\begin{lstlisting}[language=bash, numbers=none]
# Requires Emscripten SDK installed and activated
emmake make
\end{lstlisting}

The \lstinline|platform()| built-in returns \lstinline|"wasm"| when running in this mode.

\subsection{Verifying the Installation}\label{subsec:verify-install}

Once you have built \texttt{clat}, verify it works:

\begin{lstlisting}[language=bash, numbers=none]
./clat --version
\end{lstlisting}

You should see:

\begin{lstlisting}[language=bash, numbers=none]
Lattice v0.3.28
\end{lstlisting}

If you want the full help screen with all available subcommands and flags:

\begin{lstlisting}[language=bash, numbers=none]
./clat --help
\end{lstlisting}

%% ───────────────────────────────────────────────────────────
\section{Your First Program: Hello, World!}\label{sec:hello-world}
\index{Hello, World}

Let's write our first Lattice program.
Create a file called \filename{hello.lat}:

\begin{lstlisting}[language=Lattice, caption={hello.lat --- your first Lattice program}]
print("Hello, World!")
\end{lstlisting}

Run it:

\begin{lstlisting}[language=bash, numbers=none]
./clat hello.lat
\end{lstlisting}

Output:

\begin{lstlisting}[language=bash, numbers=none]
Hello, World!
\end{lstlisting}

That's it.
No \lstinline|main()| function required, no boilerplate, no imports.
Lattice programs are sequences of statements, and the interpreter executes them top to bottom.

\subsection{A Slightly More Interesting Program}\label{subsec:more-interesting}

Let's make it personal:

\begin{lstlisting}[language=Lattice, caption={greeting.lat --- string interpolation}]
let name = "Lattice"
let year = 2026
print("Welcome to ${name}! The year is ${year}.")
\end{lstlisting}

Output:

\begin{lstlisting}[language=bash, numbers=none]
Welcome to Lattice! The year is 2026.
\end{lstlisting}

Notice the \lstinline|${...}|\index{string interpolation} syntax inside double-quoted strings.
Any expression can go inside those braces---variable lookups, arithmetic, even function calls.
We will cover string interpolation in depth in \cref{ch:values-types}.

\subsection{Functions}\label{subsec:first-functions}
\index{functions!introduction}

Here is a program with a function:

\begin{lstlisting}[language=Lattice, caption={A function with a type annotation}]
fn fahrenheit_to_celsius(temp: Float) -> Float {
    (temp - 32.0) * 5.0 / 9.0
}

let boiling = fahrenheit_to_celsius(212.0)
print("Water boils at ${boiling} Celsius")
// Water boils at 100.0 Celsius
\end{lstlisting}

A few things to notice:

\begin{itemize}
  \item Function parameters have \textbf{type annotations}\index{type annotations}: \lstinline|temp: Float| declares that \lstinline|temp| must be a \lstinline|Float|. Pass the wrong type and Lattice will tell you---at runtime, with a clear error message.
  \item The \textbf{last expression} in a function body is the implicit return value\index{implicit return}. No \lstinline|return| keyword needed (though you can use one if you prefer explicit style).
  \item The \lstinline|-> Float| after the parameter list declares the return type.
\end{itemize}

\subsection{Running the Examples Directory}\label{subsec:running-examples}
\index{examples}

Lattice ships with an \filename{examples/} directory containing programs that demonstrate different features.
Try running the Fibonacci example:

\begin{lstlisting}[language=bash, numbers=none]
./clat examples/fibonacci.lat
\end{lstlisting}

This program computes Fibonacci numbers using both iterative and recursive approaches, verifies they produce the same results, and approximates the golden ratio.
Here is a simplified version of what it does:

\begin{lstlisting}[language=Lattice, caption={A taste of the Fibonacci example}]
fn fib_iterative(n: Int) -> Int {
    if n <= 1 { return n }
    flux a = 0
    flux b = 1
    flux i = 2
    while i <= n {
        let temp = a + b
        a = b
        b = temp
        i += 1
    }
    return b
}

for i in 0..10 {
    print("fib(${to_string(i)}) = ${to_string(fib_iterative(i))}")
}
\end{lstlisting}

Notice the \lstinline|flux|\index{flux} keyword: it declares mutable variables.
We will explain \lstinline|flux|, \lstinline|let|, and \lstinline|fix| in \cref{ch:variables-phase}.

%% ───────────────────────────────────────────────────────────
\section{Running Files vs.\ the REPL}\label{sec:files-vs-repl}
\index{REPL!introduction}

There are two main ways to run Lattice code: from a file, or interactively in the REPL\@.

\subsection{Running from a File}\label{subsec:running-file}

You've already seen this:

\begin{lstlisting}[language=bash, numbers=none]
./clat my_program.lat
\end{lstlisting}

Lattice reads the entire file, compiles it to bytecode, and executes it on the stack VM\@.
If you want to use the tree-walking interpreter instead (slower, but useful for debugging), pass \lstinline|--tree-walk|:

\begin{lstlisting}[language=bash, numbers=none]
./clat --tree-walk my_program.lat
\end{lstlisting}

You can also \emph{pre-compile} a file to bytecode\index{bytecode compilation} and run the compiled output later:

\begin{lstlisting}[language=bash, numbers=none]
./clat compile program.lat -o program.latc
./clat program.latc
\end{lstlisting}

The \filename{.latc} file contains serialized bytecode.
Running a pre-compiled file skips the lexing, parsing, and compilation steps, which can speed up startup for larger programs.

\begin{notebox}{Three Execution Backends}
Lattice has three ways to execute your code: the \textbf{stack-based bytecode VM} (the default and fastest), the \textbf{register-based VM} (via \lstinline|--regvm|), and the \textbf{tree-walking interpreter} (via \lstinline|--tree-walk|).
All three support the full language.
We will cover the backends in depth in Chapter~33 (\emph{The Three Backends})---for now, the default is the right choice.
\end{notebox}

\subsection{The Interactive REPL}\label{subsec:interactive-repl}

Launch the REPL by running \texttt{clat} with no arguments:

\begin{lstlisting}[language=bash, numbers=none]
./clat
\end{lstlisting}

You'll see a welcome banner:

\begin{lstlisting}[language=bash, numbers=none]
Lattice v0.3.28 -- crystallization-based programming language
Copyright (c) 2026 Alex Jokela. BSD 3-Clause License.
Type expressions to evaluate. Ctrl-D to exit.
\end{lstlisting}

Now you can type expressions and see results immediately:

\begin{lstlisting}[language=Lattice, numbers=none]
lattice> 2 + 2
=> 4
lattice> "hello" + " " + "world"
=> "hello world"
lattice> [1, 2, 3].map(|x| x * 10)
=> [10, 20, 30]
\end{lstlisting}

The \lstinline|=>|\index{=> prefix (REPL)} prefix is how the REPL shows you the result of an expression.
Statements that produce no meaningful value---like variable assignments---are silently suppressed:

\begin{lstlisting}[language=Lattice, numbers=none]
lattice> let name = "Lattice"
lattice> print(name)
Lattice
lattice>
\end{lstlisting}

The REPL maintains \textbf{persistent state}\index{REPL!persistent state}: variables, functions, structs, and enums you define in one line are available on the next.
This makes it a wonderful tool for exploratory programming.

\begin{lstlisting}[language=Lattice, numbers=none]
lattice> fn square(x: Int) -> Int { x * x }
lattice> square(7)
=> 49
lattice> let values = [1, 2, 3, 4, 5]
lattice> values.map(|x| square(x))
=> [1, 4, 9, 16, 25]
\end{lstlisting}

We will dive much deeper into the REPL in \cref{ch:repl}.

\subsection{Passing Arguments to Scripts}\label{subsec:script-args}
\index{command-line arguments}\index{args()}

Arguments after the filename are passed to your program and can be retrieved with the \lstinline|args()|\index{args()} built-in:

\begin{lstlisting}[language=Lattice, caption={echo.lat --- reading command-line arguments}]
let arguments = args()
print("You passed ${to_string(arguments.len())} arguments:")
for arg in arguments {
    print("  ${arg}")
}
\end{lstlisting}

\begin{lstlisting}[language=bash, numbers=none]
./clat echo.lat hello world 42
\end{lstlisting}

Output:

\begin{lstlisting}[language=bash, numbers=none]
You passed 4 arguments:
  echo.lat
  hello
  world
  42
\end{lstlisting}

The first element is the script's own filename, following the convention of most languages.

%% ───────────────────────────────────────────────────────────
\section{A Taste of What Makes Lattice Different}\label{sec:taste-different}
\index{Lattice!unique features}

Before we close this chapter, let's take a quick tour of some features that set Lattice apart.
Don't worry about understanding every detail yet---each of these gets a full chapter later.
The goal here is to whet your appetite.

\subsection{The Phase System in Action}\label{subsec:phase-taste}
\index{phase system!first example}

Here is the phase system at work in a realistic scenario.
Imagine building a server configuration that should be mutable during setup, then locked down before the server starts:

\begin{lstlisting}[language=Lattice, caption={Building and freezing a configuration}]
// Build the config in a fluid state
flux config = Map::new()
config["host"] = "0.0.0.0"
config["port"] = 8080
config["max_connections"] = 256

// Lock it down
freeze(config)

// Now config is crystal --- safe to share, impossible to change
print(config["host"])           // 0.0.0.0
print(phase_of(config))        // crystal

// config["port"] = 9090      -- this would error!
\end{lstlisting}

The \lstinline|forge|\index{forge} block makes this pattern even more elegant:

\begin{lstlisting}[language=Lattice, caption={Forge: build fluid, end crystal}]
fix config = forge {
    flux temp = Map::new()
    temp.set("host", "0.0.0.0")
    temp.set("port", "8080")
    temp.set("debug", "false")
    freeze(temp)
}
// config is now crystal, built in one clean block
print(config.get("host"))  // 0.0.0.0
\end{lstlisting}

\subsection{Pattern Matching}\label{subsec:match-taste}
\index{pattern matching!first example}

Lattice's \lstinline|match| expression supports literal patterns, ranges, guards, destructuring, and the compiler warns you if your patterns are not exhaustive:

\begin{lstlisting}[language=Lattice, caption={Pattern matching}]
fn classify_temperature(temp: Int) -> String {
    match temp {
        x if x < 0 => "freezing",
        0..15 => "cold",
        15..25 => "comfortable",
        25..35 => "warm",
        _ => "hot"
    }
}

print(classify_temperature(22))   // comfortable
print(classify_temperature(-5))   // freezing
print(classify_temperature(38))   // hot
\end{lstlisting}

\subsection{Structured Concurrency}\label{subsec:concurrency-taste}
\index{concurrency!first example}

Lattice uses \lstinline|scope|\index{scope} and \lstinline|spawn|\index{spawn} for structured concurrency.
Every spawned task must complete before its enclosing scope exits, and values sent across threads must be \emph{frozen}---the phase system ensures data-race safety at the language level:

\begin{lstlisting}[language=Lattice, caption={Structured concurrency with channels}]
let results = Channel::new()

scope {
    spawn {
        results.send(freeze("task A complete"))
    }
    spawn {
        results.send(freeze("task B complete"))
    }
}

print(results.recv())   // task A complete (or B, depending on scheduling)
print(results.recv())   // the other one
\end{lstlisting}

Notice: only \lstinline|freeze()|\index{freeze}'d values can be sent on a channel.
This is the phase system doing its job---guaranteeing that shared data is immutable.

\subsection{Closures and Functional Style}\label{subsec:closures-taste}
\index{closures!first example}

Closures are first-class citizens in Lattice:

\begin{lstlisting}[language=Lattice, caption={Closures and higher-order functions}]
let numbers = [1, 2, 3, 4, 5, 6, 7, 8, 9, 10]

let evens = numbers.filter(|n| n % 2 == 0)
let doubled = evens.map(|n| n * 2)
let total = doubled.reduce(|acc, n| acc + n, 0)

print(total)  // 60
\end{lstlisting}

And you can chain these together more fluidly with the pipe operator\index{pipe operator}:

\begin{lstlisting}[language=Lattice, caption={Chaining transforms with iterators}]
let total = iter([1, 2, 3, 4, 5, 6, 7, 8, 9, 10])
    .filter(|n| n % 2 == 0)
    .map(|n| n * 2)
    .reduce(0, |acc, n| acc + n)

print(total)  // 60
\end{lstlisting}

\subsection{String Interpolation Everywhere}\label{subsec:interpolation-taste}
\index{string interpolation!first example}

Double-quoted strings support \lstinline|${expr}| interpolation.
Single-quoted strings are raw---no interpolation, no escape sequences.
Triple-quoted strings preserve whitespace and support interpolation:

\begin{lstlisting}[language=Lattice, caption={String varieties}]
let language = "Lattice"
let version = version()

// Double-quoted: interpolation
print("Running ${language} ${version}")

// Single-quoted: raw, no interpolation
print('The syntax is ${expr}')  // literal ${expr}

// Triple-quoted: multi-line
let banner = """
    =============================
      Welcome to ${language}
      Version: ${version}
    =============================
    """
print(banner)
\end{lstlisting}

\subsection{One More Thing: Strict Mode}\label{subsec:strict-taste}
\index{strict mode!introduction}

At the top of any file, you can write \lstinline|#mode strict|\index{\#mode strict} to opt into stricter phase enforcement.
In strict mode, the phase checker runs before execution and catches potential issues early:

\begin{lstlisting}[language=Lattice, caption={Strict mode}]
#mode strict

flux data = [1, 2, 3]
fix frozen = freeze(data)
// In strict mode, `data` is consumed by freeze().
// Accessing `data` after this line is a compile-time error.
\end{lstlisting}

Strict mode is entirely opt-in.
The default ``casual'' mode is more forgiving, making it friendlier for scripting and exploration.
We will compare the two modes in \cref{ch:variables-phase}.

%% ───────────────────────────────────────────────────────────
\section{Exercises}\label{sec:ch01-exercises}

\begin{enumerate}
  \item \textbf{Hello, You.}
  Write a program that asks for the user's name using \lstinline|input("What is your name? ")| and prints a personalized greeting using string interpolation.

  \item \textbf{Temperature Table.}
  Write a program that prints a conversion table from Fahrenheit to Celsius for every 10 degrees from 0 to 212.
  Use a \lstinline|for| loop and the range operator.
  Hint: \lstinline|for f in range(0, 220, 10) { ... }|

  \item \textbf{Explore the REPL\@.}
  Launch the REPL (\lstinline|./clat|) and experiment.
  Define a function, call it with different arguments, and try to break it by passing the wrong types.
  What error messages do you get?

  \item \textbf{Freeze Tag.}
  In the REPL, create an array with \lstinline|flux items = [10, 20, 30]|.
  Push a new value onto it.
  Then freeze it with \lstinline|freeze(items)| and try to push again.
  Read the error message---it will become familiar.

  \item \textbf{Read an Example.}
  Read the source of \filename{examples/phase\_demo.lat} and run it.
  Before running, predict what each \lstinline|print| statement will output.
  Were you right?
\end{enumerate}

%% ───────────────────────────────────────────────────────────
\section*{What's Next}\label{sec:ch01-whats-next}

You have installed Lattice, written your first program, and seen glimpses of the phase system, pattern matching, and concurrency.
In the next chapter, we will spend quality time with the REPL---Lattice's interactive laboratory.
The REPL is where you will do most of your exploring, testing ideas, and building intuition for how the language works.
Let's set up your lab.
